\section{Erfüllung der Zielsetzung}
Ziel dieser Arbeit war die Konzeptionierung, Entwicklung und Umsetzung einer WordClock mit informativen Zusatzfunktionen.  
Hierzu zählten die Integration externer Wetter- und Sensordaten.  
Darüber hinaus sollten eine netzwerkbasierte Steuerung sowie eine modulare Hardware- und Softwarearchitektur umgesetzt werden.  

Die entwickelte WordClock erfüllt die grundlegenden Zielsetzungen der Arbeit.  
Die sprachbasierte Zeitdarstellung konnte um eine sekundengenaue Anzeige sowie zusätzliche Informationsfunktionen erweitert werden.  
Dadurch entstand ein Gesamtsystem, das über die Funktion klassischer WordClocks hinausgeht und zusätzliche Umgebungsinformationen integriert.   
Die Kombination aus Wetterdaten, Temperaturinformationen und lokaler Sensorik erweitert den Funktionsumfang der WordClock zu einem multifunktionalen Informationssystem.  
Zusätzlich ermöglicht die implementierte Weboberfläche eine flexible Benutzerinteraktion sowie die Anpassung verschiedener Anzeigeparameter.  

Die gewählte modulare Struktur erleichterte die Entwicklung und Optimierung einzelner Funktionen und bietet eine Grundlage für die Weiterentwicklung.  
Dadurch konnten Hardware- und Softwarekomponenten unabhängig voneinander entwickelt, getestet und erweitert werden.  
Insbesondere während der Fehlersuche und Integration erwies sich diese Struktur als vorteilhaft.  
Auch die mechanische Konstruktion konnte entsprechend der Zielsetzung umgesetzt werden.  
 
Die entwickelte WordClock stellt einen funktionsfähigen und erweiterbaren Prototyp mit zusätzlichen Informations- und Steuerungsfunktionen dar und erfüllt die definierten Ziele.